% =====================================================================
% PT_RH_HP_CANONICAL — Reformulation canonique de Hilbert-Polya via PT
% Auteur : Yan Senez
% Date : 2026-05-16
% Version : v5 (Phases A + B + C complétées, annexe Kähler-Fisher incluse)
%
% Source : draft_fr.md (markdown authoritative)
% Conversion en LaTeX pour soumission.
% =====================================================================

\documentclass[11pt, a4paper]{article}

% Préambule commun PT français
\usepackage{../shared/pt-common-fr}

% Spécifiques à cet article
\usepackage{geometry}
\geometry{margin=2.5cm}

% Numérotation des équations par section
\numberwithin{equation}{section}

% Métadonnées
\title{Reformulation canonique de Hilbert-Polya \\
       via la Théorie de la Persistance}
\author{Yan Senez}
\date{\today}

% =====================================================================
\begin{document}

\maketitle

\begin{abstract}
L'apport principal de cet article est de \textbf{dériver toutes les constantes
structurelles du programme Hilbert-Polya depuis un unique théorème
d'arithmétique élémentaire}. Plus précisément, on établit la chaîne
\[
  \underbrace{\text{T1 (transitions interdites mod 3)}}_{[\mathrm{THM}]\text{ classique}}
  \xrightarrow{\text{Th. 5.2}}
  \underbrace{c_1(\mathbb{S}_+) = s = 1/2}_{[\mathrm{THM}]\text{ mod Kähler-Fisher}}
  \xrightarrow{\text{Th. 3.13}}
  \underbrace{\theta_{\mathrm{PT}} = \pi}_{[\mathrm{THM}]}
\]
qui détermine canoniquement l'opérateur de Berry-Keating-PT
$\tilde H_{\mathrm{PT}} = -i\partial_\xi$ sur $L^2([0, r], d\xi)$ avec
condition aux bords antipériodique forcée par la cohomologie spinorielle.

\textbf{Théorème principal} : Inconditionnellement sous T1 et le théorème
de Kähler-Fisher (Annexe A), la conjecture de Riemann est équivalente à
la propriété spectrale : la fonction $T(s) = -d\log R/ds$ (où
$R(s) = \zeta(s)/(\zeta_+(s)\zeta_-(s))$) est holomorphe sur
$\{0 < \mathrm{Re}(s) < 1,\; \mathrm{Re}(s) \neq 1/2\}$.

L'équivalence est rigoureuse et \emph{inconditionnelle}. Elle ne prouve pas RH —
elle le \emph{reformule canoniquement} dans un cadre où tous les ingrédients
du programme Hilbert-Polya sont dérivés sans constante ad hoc.

À la connaissance de l'auteur, c'est la première formulation du programme
Hilbert-Polya où aucune constante n'est postulée \emph{ad hoc}, et où RH
apparaît comme propriété spectrale d'un opérateur canoniquement déterminé.
\end{abstract}

\tableofcontents
\bigskip
\hrule
\bigskip

% =====================================================================
\section{Introduction}
\label{sec:intro}

\subsection{Le programme Hilbert-Polya}

L'observation de Hilbert et Pólya, vers 1920, que les zéros non triviaux de la
fonction zêta de Riemann pourraient être les valeurs propres d'un opérateur
auto-adjoint sur un espace de Hilbert reste, plus d'un siècle après, la
direction la plus suggestive vers une démonstration spectrale de l'hypothèse
de Riemann (RH).

Le formalisme a reçu deux avancées majeures. Berry et Keating \cite{BerryKeating1999},
en 1999, ont proposé un candidat semi-classique explicite : l'opérateur de
dilatation $H = (xp + px)/2$ sur la demi-droite. Connes \cite{Connes1999}
a généralisé ce cadre à un opérateur de scaling sur l'espace adélique des
classes d'idèles. Ces deux approches laissent plusieurs ingrédients
\textbf{postulés ad hoc} : cellule de Planck $2\pi$, régularisation, condition
aux bords, scattering matrix, résidu cohomologique $1/8$, spin $s = 1/2$.

% TODO: Convertir §1.2-§1.4 du draft_fr.md

% =====================================================================
\section{Prérequis classiques}
\label{sec:prerequis}

% TODO: Convertir §2.1, §2.2, §2.3 du draft_fr.md
% Théorèmes : 2.2 (T1), 2.4 ($s=1/2$), 2.6 (non-annulation), 2.8 (Lemme bornes), 2.10 (Fredholm)

% =====================================================================
\section{L'opérateur canonique $H_{\mathrm{PT\text{-}BK}}$}
\label{sec:operateur}

% TODO: Convertir §3.1 à §3.6 du draft_fr.md
% Théorèmes : 3.2 (indices), 3.6 (spectre), 3.9 (anticommutation), 3.10 (sélection BC), 3.13 (Hyp 3.11 promu)

% =====================================================================
\section{Quatre mécanismes équivalents pour $\mathrm{Re}(s) = 1/2$}
\label{sec:mecanismes}

% TODO: Convertir §4.1 à §4.5 du draft_fr.md

% =====================================================================
\section{Sur-détermination cohomologique du résidu $1/8$}
\label{sec:cohomologie}

% TODO: Convertir §5.1 à §5.5 du draft_fr.md
% Théorèmes : 5.2 (c_1 = 1/2), 5.6 (sur-détermination)

% =====================================================================
\section{Construction du résidu et de la trace régularisée}
\label{sec:trace}

% TODO: Convertir §6.1 à §6.6 du draft_fr.md
% Théorèmes : 6.2 (formule de trace), 6.7 (existence prolongement distributionnel)

% =====================================================================
\section{Reformulation rigoureuse de RH}
\label{sec:rh-reformulation}

% TODO: Convertir §7.1 à §7.5 du draft_fr.md
% Théorème principal : 7.1

% =====================================================================
\section{Comparaison avec Berry-Keating et Connes}
\label{sec:comparaison}

% TODO: Convertir §8.1 à §8.4 du draft_fr.md

% =====================================================================
\section{Validation numérique sommaire}
\label{sec:validation}

Détails complets dans l'article séparé \texttt{PT\_RH\_VALIDATION/}.
Résumé :
\begin{itemize}
  \item 75/79 zéros de $\zeta$ captés à $|\Delta| < 0.5$ sur $t \in [10, 200]$
        (régularisation R1, $\varepsilon = 0.2$)
  \item 267/269 zéros captés à $|\Delta| < 1.5$ sur $t \in [10, 500]$ (99.3\,\%)
  \item Test de chance Weyl rejeté à $25\sigma$ sur les 13 premiers zéros
\end{itemize}

% =====================================================================
\section{Limites et ouvertures}
\label{sec:limites}

% TODO: Convertir §10.1 à §10.4 du draft_fr.md

% =====================================================================
\section{Conclusion}
\label{sec:conclusion}

% TODO: Convertir §11 du draft_fr.md

% =====================================================================
\appendix

\section{Théorème de Kähler-Fisher}
\label{app:kahler-fisher}

% TODO: Convertir Annexe A du draft_fr.md (A.1 à A.6)
% Théorème A.6 : Kähler-Fisher complet

% =====================================================================
\bibliographystyle{alpha}
\begin{thebibliography}{99}

\bibitem{BerryKeating1999}
M.~V. Berry and J.~P. Keating.
\newblock $H = xp$ and the Riemann zeros.
\newblock In \emph{Supersymmetry and Trace Formulae: Chaos and Disorder},
NATO ASI Series B 370, Plenum Press, 1999, pp. 355--367.

\bibitem{Connes1999}
A.~Connes.
\newblock Trace formula in noncommutative geometry and the zeros of the
Riemann zeta function.
\newblock \emph{Selecta Mathematica}, 5:29--106, 1999.

\bibitem{Riemann1859}
B.~Riemann.
\newblock Über die Anzahl der Primzahlen unter einer gegebenen Grösse.
\newblock \emph{Monatsberichte der Berliner Akademie}, 671--680, 1859.

\bibitem{Titchmarsh1986}
E.~C. Titchmarsh.
\newblock \emph{The Theory of the Riemann Zeta-Function}.
\newblock 2nd edition, ed. D.~R. Heath-Brown, Oxford University Press, 1986.

\bibitem{Landau1909}
E.~Landau.
\newblock \emph{Handbuch der Lehre von der Verteilung der Primzahlen}.
\newblock Teubner, Leipzig, 1909.

\bibitem{LandauWalfisz1920}
E.~Landau and A.~Walfisz.
\newblock Über die Nichtfortsetzbarkeit einiger durch Dirichletsche Reihen
definierter Funktionen.
\newblock \emph{Rendiconti del Circolo Matematico di Palermo}, 44:82--86, 1920.

\bibitem{Cohen2007}
H.~Cohen.
\newblock \emph{Number Theory Vol. II: Analytic and Modern Tools}.
\newblock Springer GTM 240, 2007.

\bibitem{ReedSimon1980}
M.~Reed and B.~Simon.
\newblock \emph{Methods of Modern Mathematical Physics II: Fourier Analysis,
Self-Adjointness}.
\newblock Academic Press, 1980.

\bibitem{AtiyahSinger1968}
M.~F. Atiyah and I.~M. Singer.
\newblock The index of elliptic operators on compact manifolds.
\newblock \emph{Bulletin of the AMS}, 69:422--433, 1968.

\bibitem{APS1975}
M.~F. Atiyah, V.~K. Patodi, and I.~M. Singer.
\newblock Spectral asymmetry and Riemannian geometry.
\newblock \emph{Mathematical Proceedings of the Cambridge Philosophical Society},
77:43--69 (I); 78:405--432 (II); 79:71--99 (III), 1975.

\bibitem{Berry1984}
M.~V. Berry.
\newblock Quantal phase factors accompanying adiabatic changes.
\newblock \emph{Proceedings of the Royal Society A}, 392:45--57, 1984.

\bibitem{Hadamard1896}
J.~Hadamard.
\newblock Sur la distribution des zéros de la fonction $\zeta(s)$ et ses
conséquences arithmétiques.
\newblock \emph{Bulletin de la SMF}, 24:199--220, 1896.

\bibitem{Cencov1982}
N.~N. Čencov.
\newblock \emph{Statistical Decision Rules and Optimal Inference}.
\newblock AMS Translations of Mathematical Monographs 53, 1982.

\bibitem{NewlanderNirenberg1957}
A.~Newlander and L.~Nirenberg.
\newblock Complex analytic coordinates in almost complex manifolds.
\newblock \emph{Annals of Mathematics}, 65:391--404, 1957.

\bibitem{GriffithsHarris}
P.~Griffiths and J.~Harris.
\newblock \emph{Principles of Algebraic Geometry}.
\newblock Wiley, 1978.

\bibitem{LMFDB2024}
LMFDB Collaboration.
\newblock \emph{The L-functions and Modular Forms Database}.
\newblock \url{https://www.lmfdb.org/}, 2024.

\bibitem{SenezMonographie2026}
Y.~Senez.
\newblock \emph{Théorie de la Persistance : du paramètre arithmétique
$s = 1/2$ aux observables du Modèle Standard}.
\newblock Monographie en cours, 2026.

\end{thebibliography}

\end{document}
